\section{Consequences and remaining questions}\label{sec:discussion}

\subsection{Relaxation, mixing, and the visible modes}

For the polymer semigroup $P_t=\exp(tL_M)$, the exact reduction gives
\begin{equation}\label{eq:L2-decay}
 \|P_tf\|_{L^2(\pi_M)}\le e^{-\gamma_Qt}\|f\|_{L^2(\pi_M)},
 \qquad \pi_Mf=0.
\end{equation}
The exponent is optimal. A strictly monotone first eigenfunction
of the count chain attains it. At stationarity,
\begin{equation}\label{eq:cov-decay}
 \left|\operatorname{Cov}_{\pi_M}(f(M_0),f(M_t))\right|
 \le e^{-\gamma_Qt}\operatorname{Var}_{\pi_M}(f).
\end{equation}
For $\beta=b/n$, relaxation therefore occurs on the scale
$\sqrt n/(2\sqrt b)$, while total-variation mixing from either
extreme polymer configuration has the larger scale
$\sqrt n\log n/(8\sqrt b)$. The logarithm measures the time needed
for the count to enter its equilibrium fluctuation range.
Its equilibrium location is of order $\sqrt n$, its standard
deviation is of order $n^{1/4}$, and its local restoring rate is
asymptotic to $2\sqrt b/\sqrt n$.

The distinction between labelled and unlabelled states is necessary
when fragmentation is rare. At $n=2$, the polymer chain has only
one dimer and two monomers, with gap $\beta+1$. The labelled cut
chain has an additional eigenfunction given by the difference of
the two singleton indicators, with eigenvalue $-\beta$.
Both singleton cuts represent the same dimer, so this function
cannot be a polymer observable. In general, rotation invariance
removes degree one, and the sector comparison proves that no
remaining polymer mode is slower than the count mode.

With spatial diffusion the labels are part of the state. The
degree-one mode is then visible, and the exact minimum
$\min\{\gamma_Q,\eta_1+D\lambda_G\}$ determines which mode controls
relaxation. At $\beta=b/n$ its two reaction rates are asymptotic to
$2\sqrt b/\sqrt n$ and $\sqrt b/\sqrt n$. At
$\beta=b/n^2$ the corresponding rates are $\rho(b)/n$ and
$\sigma(b)/n$. These formulas give both the crossover and the
diffusion strength needed for count relaxation to become the
limiting mechanism.

\subsection{The remaining mixing question}

The count gap and the polymer gap agree, but a spectral-gap
identity does not by itself identify the worst-case mixing time;
see~\cite{LevinPeres2017}. From a specified state $m$, the usual
estimate is
\[
 \|P_t(m,\cdot)-\pi_M\|_{\mathrm{TV}}
 \le\tfrac12\sqrt{\pi_M(m)^{-1}-1}\,e^{-\gamma_Qt}.
\]
Its prefactor can be large. The singleton-fiber argument avoids
that prefactor for the two extreme states, because their evolving
density is exactly constant on count fibers. An arbitrary initial
partition has additional within-fiber information.
Theorem~\ref{thm:full-polymer-tv-upper-new} controls that information
uniformly and yields, for each fixed $\varepsilon\in(0,1)$,
\[
 \frac{\sqrt n\log n}{8\sqrt b}-O_{b,\varepsilon}(\sqrt n)
 \le t^{\mathrm{polymer}}_{\mathrm{mix},n}(\varepsilon)
 \le C_b\sqrt n(\log n)^2.
\]

The main open mixing problem is to determine whether this
additional information disappears by
$\sqrt n\log n/(8\sqrt b)+O_b(\sqrt n)$ uniformly over initial
partitions, or whether some configurations mix later. The
endpoint cutoff theorem supplies the lower bound with the exact
constant. It does not identify the worst initial polymer state,
prove a worst-case polymer cutoff, or determine a limiting cutoff
profile. An optimal window lower bound is also a separate question.

\subsection{Beyond the solvable reaction structure}

The bounded-conductance comparisons allow spatial and size
dependence while retaining a specified equilibrium law. Growing
kernels require further estimates because a uniform comparison
with the constant kernel loses powers of $n$.
For the shared multipliers $ij$ and $i+j$, the monomer test function
gives upper bounds of order $1+\beta$
(Proposition~\ref{prop:growing-kernel-upper}). At $\beta=b/n$, the
proved lower bounds still have order $n^{-1/2}$, leaving a gap between
the two estimates. The classical Ewens
model shows that the fragmentation rule matters as much as the
merger kernel: its equilibrium weights and its spectral scale are
different from those of the constant-kernel polymer chain.

The nonreversible results distinguish two convergence quantities.
A directed current can improve the least nonzero real part of the
spectrum while leaving the symmetric part, and hence the best
prefactor-one $L^2$ exponent, unchanged. The four-state polymer
construction proves strict acceleration for each finite system.
It leaves open a choice of currents with a gain bounded away from
zero as $n$ grows, as well as a quantitative total-variation
comparison at the cutoff scale.

Finally, the spatial model studied here has local conservative
motion but globally count-dependent reactions. A general spatial
coagulation--fragmentation system may have spatially varying
concentrations, cluster diffusion depending on size, collision
kernels depending on distance, and a nonexchangeable equilibrium.
Those changes remove the tensor decomposition used in the graph
formula. Establishing comparable estimates for such systems would
require new control of their spatial and cluster-size fluctuations.
